% ------------------------------------------------------------- Conclusion
A UAV's navigation stack cannot tell an attacker from the weather, and an
attacker who knows this will dress accordingly. We turned that observation into
a threat model, a requirement and a measurement. The threat model is
$\varepsilon$-observational equivalence, which bounds every detector's advantage
by a total-variation distance. The requirement is that a mission-level
guarantee must bound undetected failure without relying on the detector to
notice the disguise. The measurement is a cross-layer benchmark with enforced
provenance and a cross-evaluated evasion study. In that study every disguise
evaded every detector, and the deviation a detector admitted depended on
whether it modelled the disguise as benign, unless a trusted re-anchoring
alarm capped it for all of them.

The unified budget meets the requirement. Its stochastic terms are aggregated
over every coupling, including one an attacker imposes by controlling when
jamming happens; its deterministic terms add no dependence gap; and its single
empirical adversarial input is reported together with the case in which an
independent attacker exceeded it. The analysis also corrects a hypothesis from
our proposal: the union bound is not where the looseness lives. The slack lies
in the dependence, envelope and temporal gaps. In our pilot the temporal gap
was the largest at moderate budgets, and it is the most tractable. For the defender, re-anchoring and closing the
terminal interval mattered more than any threshold. The next step is to
re-establish these results in PX4 software-in-the-loop simulation and with
measured marginals, as the released runbook specifies.
